\documentclass[12pt]{article}
\usepackage[utf8]{inputenc}
\usepackage[T1]{fontenc}
\usepackage{hyperref}
\title{Mayi-Yapi Language Resources for K-12 Curriculum Development}
\author{Compiled by Claude Sonnet 5.5 (claude-sonnet-5-5) for Andr\'as Kornai}
\date{}
\begin{document}
\maketitle

\section{Introduction}
Mayi-Yapi (ISO 639-3: xyj; AIATSIS code G20) is an extinct Australian language of the Mayi group, from the Gulf Country of Queensland. Searches found very little material beyond catalogue entries and one scholarly survey of the group. No teaching materials, native-language primers or revitalization programmes turned up. This report lists only what the searches showed.

\section{Primary Educational Resources}
No primers, textbooks or school programmes for Mayi-Yapi were found.

\section{Linguistic and Reference Materials}
\begin{itemize}
\item AIATSIS Austlang entry G20, Mayi-Yapi. It gives the traditional location (Cloncurry River region, Flinders River, Leichhardt-Alexandra Divide) and notes that \emph{mayi} means ``speech, language''. \url{https://collection.aiatsis.gov.au/austlang/language/g20}
\item Gavan Breen (1981), \emph{The Mayi Languages of the Queensland Gulf Country}, AIAS New Series 29. This is the principal reference for the group. Only the Glottolog bibliographic record was seen. \url{https://glottolog.org/resource/reference/id/64840}
\item Glottolog language record for the ISO code mnt, which groups several Mayi varieties. It notes that some authors treat the Mayi-Yapi people as distinct from their neighbours. \url{https://glottolog.org/resource/languoid/iso/mnt}
\end{itemize}

\section{Religious and Community Resources}
No community organisation or religious materials specific to Mayi-Yapi were found in the searches.

\section{Digital and Audio Resources}
No audio, video or software resources were found.

\section{Curriculum Development Context}
The language has no living speakers and very sparse documentation, and the other Mayi languages are in the same state. A curriculum would have to rest on Breen's 1981 work and on archival holdings, which can be requested through AIATSIS. A curriculum should be developed only with the descendant community, and the cultural content should come from them. Mayi-Yapi has no identified community programme to point to.
\end{document}